\documentclass{article}
\usepackage{../header}
\title{Algebraic Topology}
\author{Notes made by Finley Cooper}
\newcommand{\rel}{\mathrm{rel}\:}
\begin{document}
  \maketitle
  \newpage
  \tableofcontents
  \newpage
  \section{Introduction}
  \subsection{Recap of IB Topological Spaces}
  We use algebra to turn topological questions to algebraic questions. This results in a loss of information but lets us prove some results.
  \begin{example}
    \[
	    \R^n \cong \R^m \implies n = m
    \]
  \end{example}
  This is a statement which is very hard to prove otherwise. Likewise we can distinguish topological surfaces with a different number of `holes' in their surface.
  \par
  For this section all functions are continuous until mentioned otherwise.
  \begin{definition}
	  (Map) A \textit{map} is a continuous function.
  \end{definition}
  \begin{lemma}
	  (Gluing lemma) Let $ f:X\to Y $ be a function and $ C, K\subset X$ be closed subsets with $ C\cup K = X $. Then $ f $ is continuous if and only if the restriction of $ f $ on $ C $ and $ K $ is continuous.
  \end{lemma}
  \pf We have that $ f $ is continuous if and only if for all closed sets $ D\subset Y $ the preimage of $ D $ is closed. Since $ \inv f(D) = (\inv f(D)\cap C)\cup(\inv f(D)\cap K) $, this holds if and only if $ \forall D \subset Y $ closed we have that $ \inv f(D) \cap C $ and $ \inv f(D)\cap K $ are closed. Since $ C $ and $ K $ are closed in $ X $, a subset of $ C $ (resp.\ $ K $) is closed in $ C $ (resp.\ $ K $) if and only if it is closed in $ X $. Hence this holds if and only if $ \inv{(f|_C)}(D), \inv{(f|_K)}(D) $ are closed in $ C $ and $ K $ respectively. Hence this holds if and only if the restriction of $ f $ on $ C $ and $ K $ is continuous.\qed
  \begin{lemma}
	  (Lebesgue number lemma) Let $ (X,d) $ be a compact metric space. For any open cover $ \mathcal U = \{U_\alpha\}_{\alpha\in I} $ of $ X $ there exists a $ \delta > 0 $ such that for all $ x\in X $, $ B_\delta(x)\subset U_\alpha $ for some $ \alpha\in I $.
  \end{lemma}
  \pf Suppose no such $ \delta $ existed. Then for all $ n > 0 $ there is some $ x_n\in X $ such that $ B_{1/n}(x_n)\not\subset U_\alpha $ for all $ \alpha\in I $. The sequence $ \{x_n\} $ admits a convergent subsequence due to compactness, so \textit{wlog} replace $ x_n $ with its convergent subsequence, and take $ x_n\to y $. Then take $ \alpha\in I $ such that $ y\in U_\alpha $. Then since $ U_\alpha $ is open there exists a $ r>0 $ such that $ B_r(y)\subset U_\alpha $. Take $ N>0 $ such that
\[
  \begin{cases}
    \frac 1N < \frac r2 \\
    d(x_N,y) < \frac r2
  \end{cases}
\]
which implies that $ B_{1/N}(x_N)\subset B_r(y)\subset U_\alpha $ which is a contradiction.\qed
\subsection{Cell complexes}
\begin{definition}
	(Cell attachment) Given a space $ X $ and a map $ f:S^{n-1} \to X $. The space $ X\cup_f D^n := X \sqcup D^n / \sim $ where $ x\sim f(x), \forall x \in S^{n-1}\subset D^n $ is the result of attaching an $ n $-cell to $ X $ along $ f $.
	\par
	The image of $ f $ is called the \textit{attaching sphere}.
\end{definition}
\begin{definition}
	(Cell complex) A \textit{finite-dimensional cell complex} is obtained by:
	\begin{enumerate}
		\item Start with a set $ X^0 $ with the discrete topology.
		\item Inductively once $ X^{n-1} $ is defined we form $ X^n $ by attaching a collection of $ n $-cells along some collection of attachment maps
			\[
				\{ f_\alpha:S^{n-1} \to X^{n-1}\}_{\alpha\in I}.
			\]
		i.e.\ that $ X^n = \left( X^{n-1} \sqcup \bigsqcup_{\alpha\in I} D_\alpha^n\right) / \sim $ where we define $ x\sim f_\alpha(x) $ for $ x\in S^{n-1}_\alpha = \partial D^n_\alpha $.
	\item $ X=X^k $ for some $ k\ge 0 $. The minimal such $ k $ is called the dimension of $ X $.
	\end{enumerate}
\end{definition}
\begin{remark}
  In general we can make sense of infinite dimensional cell complexes. Suppose that we have
  \[
    X^0 \subset X^1 \subset \cdots
  \]
  such that
  \begin{enumerate}
	  \item $ X^0 $ is discrete,
	  \item $ X^n $ is formed by attaching $ n $-cells to $ X^{n-1} $.
  \end{enumerate}
  Define
  \[
	  X = \bigcup_{n\ge 0} X^n
  \]
  with the topology where $ U \subset X $ is open $ \iff U \cap X^n $ is open in $ X^n $ for all $ n $. We say that $ X $ is a CW-complex and that $ X^n $ is the $ n $-skeleton of $ X $.
 \end{remark}
 \section{Homotopy and the fundamental group}
 \subsection{Homotopy}
 \begin{definition}
	 We denote by $ I $ the closed interval $ [0,1] $.
 \end{definition}
 \begin{definition}
	 (Homotopy) Let $ f,g:X\to Y $. Then we say that a map $ H: X\times I \to Y $ is a \textit{homotopy from} $ f $ \textit{to} $ g $ if $ H(x,0) = f(x) $ and $ H(x,1)=g(x) $. If such an $ H $ exists we say that $ f $ is \textit{homotopic} to $ g $ and $ f\simeq g $.
 \end{definition}
If $ A\subset X $ and $ f|_A = g|_A $ then a homotopy $ H $ from $ f $ to $ g $ is relative to $ A $ (i.e.\ a homotopy $ \rel A $) if $ H(a,t) = f(a) = g(a) $ for all $ a\in A $ and $ t\in I $. If such an $ H $ exists we write $ f\simeq_A g $.
\begin{proposition}
  The relation homotopic $ \rel A $ defines an equivalence relation.
\end{proposition}
\pf
\begin{enumerate}
	\item $ f\simeq_A f $ trivially by $ H(x,t) = f(x) $.
	\item $ f\simeq_A g \implies g\simeq_A f $, by replacing $ t $ with $ 1-t $.
	\item Suppose that $ f\simeq_A g, g\simeq_A h $. Then let $ H_1 $ be the homotopy $ \rel A $ from $ f $ to $ g $ and $ H_2 $ for $ g $ to $ h $. Then define
		\[H(x,t) =
		  \begin{cases}
			  H_1(x,2t) & 0 \le t \le \frac 12 \\
			  H_2(x,2t-1) & \frac 12 \le t \le 1
		  \end{cases}.
		\]
\end{enumerate}
Since $ H_1(x,1) = g(x) = H_2(x,0) $, $ H $ is well-defined, and by the gluing lemma (applied to the closed sets $ X\times[0,\frac12] $ and $ X\times[\frac12,1] $) $ H $ is continuous. Also $ H(a,t) = f(a) $ for all $ a\in A $, so $ H $ is a homotopy $ \rel A $ from $ f $ to $ h $ and therefore $ f\simeq_A h $.\qed
\begin{definition}
	(Homotopy equivalence) Given a map $ f:X\to Y $ we say that $ f $ is a \textit{homotopy equivalence} if there exists a $ g: Y\to X $ such that $ f\circ g \simeq \id_Y $ and $ g \circ f \simeq \id_X $. We call $ g $ a \textit{homotopy inverse} of $ f $. If there exists a homotopy equivalence $ X\to Y $ then $ X $ is \textit{homotopy equivalent} to $ Y $ and we write that $ X\simeq Y $.
\end{definition}
\begin{definition}
	(Contractible) A space $ X $ is \textit{contractible} if $ X $ is homotopy equivalent to the one-point space $ \{*\} $.
\end{definition}
\begin{lemma}
  Let $ f_0,f_1 : X\to Y $ be homotopic maps and $ g_0,g_1 : Y\to Z $ also be homotopic maps. Then
  \[
    g_0 \circ f_0 \simeq g_1 \circ f_1.
  \]
\end{lemma}
\pf We want to show that $ g_0\circ f_0 \simeq g_0\circ f_1 \simeq g_1 \circ f_1 $.\par
Choose homotopies $ H $ from $ f_0 $ to $ f_1 $ and $ H' $ from $ g_0 $ to $ g_1 $. Then $ g_0 \circ H $ is a homotopy from $ g_0f_0 $ to $ g_0f_1 $. We also have that $ H'\circ (f_1 \times \id_I) $ is a homotopy from $ g_0f_1 $ to $ g_1f_1 $ hence we're done by transitivity.\qed
\begin{proposition}
  Homotopy equivalence is an equivalence relation.
\end{proposition}
\pf
\begin{enumerate}
	\item $ X\simeq X $ is obvious from the identity map.
	\item $ X\simeq Y $ then we have $ Y\simeq X $ just by switching the roles of $ f $ and $ g $.
	\item $ X \simeq Y $ and $ Y \simeq Z $. Then suppose $ f,g $ are the homotopy equivalence and its inverse for $ X $ and $ Y $, and $ f',g' $ for $ Y $ and $ Z $. Consider the map $ f'f $ from $ X $ to $ Z $ and its candidate homotopy inverse $ gg' $. By the lemma above,\par
		\begin{align*}
			gg'f'f \simeq g\id_Y f &= gf \simeq \id _X \\
			f'fgg'\simeq f'\id_Y g' &= f'g' \simeq \id_Z
		\end{align*}
		Therefore $ f'f $ is a homotopy equivalence and $ X\simeq Z $.\qed
\end{enumerate}
\begin{example}
  Consider the inclusion
  \[
	  \tau: S^1 \hookrightarrow \R^2\setminus \{0\}
  \]
  with homotopy inverse $ r : \R^2 \setminus \{0\} \to S^1 $ given by projecting each vector onto its unit vector, $ r(x) = x/\|x\| $. Since vector normalisation doesn't do anything to unit vectors one of the directions of composition is obvious: $ r\circ\tau = \id_{S^1} $. For the other direction,
  \[
	  \R^2 \setminus \{0\} \to S^1 \hookrightarrow \R^2 \setminus \{0\}
  \]
  we're given the homotopy
  \begin{align*}
	  H : (\R^2 \setminus \{0\}) \times I &\to \R^2\setminus \{0\}\\
	  (x,t) &\mapsto \frac{x}{t+(1-t)\|x\|}
  \end{align*}
  which is well-defined since the denominator is strictly positive, and is a homotopy from $ \tau \circ r $ to $ \id_{\R^2 \setminus \{0\}} $. Therefore we have that $ S^1 \simeq \R^2 \setminus \{0\} $. Similarly we have that $ S^{n-1} \simeq \R^n \setminus \{0\}$.
\end{example}
\begin{example}
  We have that $ \R^n $ is contractible. Let $ f:\R^n\to\{0\} $ be the constant map and $ g:\{0\}\hookrightarrow\R^n $ be the inclusion. Then
  \[
	  f\circ g = \id_{\{0\}}
  \]
  so one direction is trivial. For the other direction, $ g\circ f:\R^n\to\R^n $ is the constant map $ x\mapsto 0 $. Define the straight-line homotopy
  \begin{align*}
	  H : \R^n \times I &\to \R^n\\
	  (x,t) &\mapsto tx.
  \end{align*}
  This is continuous since scalar multiplication $ \R\times\R^n\to\R^n $ is continuous. We have $ H(x,0) = 0 = (g\circ f)(x) $ and $ H(x,1) = x = \id_{\R^n}(x) $, so $ H $ is a homotopy from $ g\circ f $ to $ \id_{\R^n} $. Therefore $ g\circ f\simeq\id_{\R^n} $, so $ f $ is a homotopy equivalence with homotopy inverse $ g $. Hence $ \R^n\simeq\{0\} $, i.e.\ $ \R^n $ is contractible.
  \par
  The same argument works for any $ X\subset\R^n $ which is star-shaped about some point $ x_0 $ (e.g.\ any convex set), using $ H(x,t) = (1-t)x_0 + tx $.
\end{example}
  \end{document}
